\textbf{Phần 1. Trắc nghiệm lựa chọn}
\setcounterpageref{ex}{0}

\begin{ex} %[4H6V30]
Để vẽ đường thẳng đi qua điểm $M$ và song song với đường thẳng $a$, ta có thể thực hiện
\choice
{\True Vẽ đường thẳng phụ vuông góc với $a$, sau đó vẽ đường thẳng qua $M$ vuông góc với đường thẳng phụ đó}
{Dùng thước kẻ một đường bất kì}
{Dùng compa quay 1 vòng tròn}
{Dùng thước đo độ đo góc $45^\circ$}
\loigiai{
Hai đường thẳng cùng vuông góc với một đường thẳng thứ ba thì song song với nhau.\\
\textbf{Đáp án của câu hỏi là A}
}
\end{ex}

\begin{ex} %[4H6V30]
Dụng cụ thường kết hợp cùng thước thẳng để vẽ hai đường thẳng song song là
\choice
{Com-pa}
{Kéo cắt giấy}
{\True Ê-ke}
{Bút chì màu}
\loigiai{
Kết hợp thước kẻ và ê-ke để trượt vẽ các đường thẳng song song.\\
\textbf{Đáp án của câu hỏi là C}
}
\end{ex}

\begin{ex} %[4H6V30]
Để vẽ một hình chữ nhật có kích thước $6\text{ cm}$ và $4\text{ cm}$, ta có thể vẽ
\choice
{3 đường thẳng song song}
{\True 2 cặp đoạn thẳng song song và vuông góc từng đôi một}
{2 đoạn thẳng cắt nhau tại trung điểm}
{1 hình tam giác vuông}
\loigiai{
Vẽ các cặp cạnh song song và vuông góc nhau.\\
\textbf{Đáp án của câu hỏi là B}
}
\end{ex}

\begin{ex} %[4H6V30]
Khi trượt một cạnh của ê-ke dọc theo mép thước kẻ thẳng, cạnh góc vuông còn lại của ê-ke sẽ vạch ra các đường thẳng
\choice
{Cắt nhau}
{Vuông góc với nhau}
{Trùng nhau}
{\True Song song với nhau}
\loigiai{
Các đường thẳng vạch ra song song với nhau.\\
\textbf{Đáp án của câu hỏi là D}
}
\end{ex}

\begin{ex} %[4H6V30]
Để vẽ đường thẳng $b$ song song với đường thẳng $a$ và cách $a$ một khoảng $3\text{ cm}$, ta lấy một điểm $M$ cách $a$ đoạn $3\text{ cm}$ rồi vẽ đường thẳng qua $M$
\choice
{\True Song song với $a$}
{Vuông góc với $a$}
{Cắt $a$ tại góc $30^\circ$}
{Cắt $a$ tại góc $45^\circ$}
\loigiai{
Vẽ đường thẳng đi qua $M$ song song với $a$.\\
\textbf{Đáp án của câu hỏi là A}
}
\end{ex}

\begin{ex} %[4H6V30]
Muốn vẽ một hình bình hành $ABCD$, sau khi vẽ hai cạnh kề $AB$ và $AD$, ta vẽ đường thẳng qua $B$ song song với $AD$ và đường thẳng qua $D$ song song với
\choice
{$AD$}
{$BD$}
{\True $AB$}
{$AC$}
\loigiai{
Qua $D$ vẽ đường thẳng song song với $AB$.\\
\textbf{Đáp án của câu hỏi là C}
}
\end{ex}

\begin{ex} %[4H6V30]
Cho 3 đường thẳng phân biệt $a, b, c$. Nếu $a // b$ và $b // c$ thì
\choice
{$a \perp c$}
{\True $a // c$}
{$a$ cắt $c$}
{$a$ trùng $c$}
\loigiai{
Tính chất bắc cầu của quan hệ song song: $a // c$.\\
\textbf{Đáp án của câu hỏi là B}
}
\end{ex}

\begin{ex} %[4H6V30]
Trên lưới ô vuông, hai đường thẳng đi qua các đường kẻ ngang cách nhau 2 ô vuông là hai đường thẳng
\choice
{Vuông góc}
{Cắt nhau}
{Trùng nhau}
{\True Song song với nhau}
\loigiai{
Các đường kẻ ngang trên lưới ô vuông luôn song song với nhau.\\
\textbf{Đáp án của câu hỏi là D}
}
\end{ex}

\begin{ex} %[4H6V30]
Để vẽ một chiếc thang gỗ có các bậc thang đều nhau, người thợ mộc cần đóng các bậc thang
\choice
{\True Song song với nhau và vuông góc với hai thanh dọc}
{Chéo nhau}
{Vuông góc với nhau}
{Tự do}
\loigiai{
Các bậc thang song song với nhau.\\
\textbf{Đáp án của câu hỏi là A}
}
\end{ex}

\begin{ex} %[4H6V30]
Cho tam giác $ABC$. Qua đỉnh $A$ vẽ đường thẳng $d$ song song với cạnh $BC$. Có thể vẽ được bao nhiêu đường thẳng $d$ như vậy?
\choice
{$0$}
{$2$}
{\True $1$ đường thẳng duy nhất}
{Vô số}
\loigiai{
Chỉ có duy nhất 1 đường thẳng qua $A$ song song với $BC$.\\
\textbf{Đáp án của câu hỏi là C}
}
\end{ex}

\begin{ex} %[4H6V30]
Gấp một tờ giấy hình chữ nhật theo chiều ngang rồi lại gấp tiếp theo nếp gấp đó ta được các nếp gấp
\choice
{Vuông góc với nhau}
{\True Song song với nhau}
{Cắt nhau tại tâm}
{Trùng nhau}
\loigiai{
Các nếp gấp cùng chiều sẽ song song với nhau.\\
\textbf{Đáp án của câu hỏi là B}
}
\end{ex}

\begin{ex} %[4H6V30]
Cho hình bình hành $ABCD$. Lấy $M$ là trung điểm $AB$, $N$ là trung điểm $CD$. Đoạn thẳng $MN$ có tính chất nào sau đây?
\choice
{Vuông góc với $AD$}
{Cắt cạnh $CD$}
{Bằng nửa $AB$}
{\True Song song với $AD$ và $BC$}
\loigiai{
$MN$ song song với hai cạnh bên $AD$ và $BC$.\\
\textbf{Đáp án của câu hỏi là D}
}
\end{ex}

\textbf{Phần 2. Trắc nghiệm đúng sai}
\setcounterpageref{ex}{0}

\begin{ex} %[4H6V30]
Xét tính đúng sai về phương pháp vẽ hình:
\choiceTF
{\True Có thể dùng ê-ke và thước thẳng để vẽ hai đường thẳng song song.}
{\True Hai đường thẳng cùng vuông góc với một đường thẳng thứ ba thì song song với nhau.}
{Qua một điểm ngoài đường thẳng ta có thể vẽ được nhiều đường thẳng song song với đường thẳng đó.}
{\True Khoảng cách giữa các vạch kẻ đường dành cho người đi bộ là các đoạn thẳng song song.}
\loigiai{
\begin{itemize}
    \item Dùng ê-ke và thước kẻ: Đúng.
    \item Cùng vuông góc thì song song: Đúng.
    \item Chỉ vẽ được duy nhất 1 đường thẳng song song (tiên đề Euclid). Sai.
    \item Vạch kẻ đường đi bộ song song: Đúng.
\end{itemize}
\begin{itemize}
    \item a) \textbf{Đúng}.
    \item b) \textbf{Đúng}.
    \item c) \textbf{Sai}.
    \item d) \textbf{Đúng}.
\end{itemize}
}
\end{ex}

\begin{ex} %[4H6V30]
Cho tam giác $ABC$. Lấy $M$ là trung điểm $AB$, vẽ đoạn $MN // BC$ ($N$ thuộc $AC$). Xét tính đúng sai:
\choiceTF
{\True Đoạn thẳng $MN$ song song với đoạn thẳng $BC$.}
{\True Điểm $N$ là trung điểm của cạnh $AC$.}
{Đoạn thẳng $MN$ vuông góc với cạnh $AB$.}
{\True Độ dài đoạn thẳng $MN$ bằng một nửa độ dài cạnh $BC$.}
\loigiai{
\begin{itemize}
    \item $MN // BC$: Đúng.
    \item $N$ là trung điểm $AC$: Đúng.
    \item $MN$ không nhất thiết vuông góc $AB$. Sai.
    \item $MN = \dfrac{1}{2} BC$ (đường trung bình): Đúng.
\end{itemize}
\begin{itemize}
    \item a) \textbf{Đúng}.
    \item b) \textbf{Đúng}.
    \item c) \textbf{Sai}.
    \item d) \textbf{Đúng}.
\end{itemize}
}
\end{ex}

\textbf{Phần 3. Trả lời ngắn}
\setcounterpageref{ex}{0}

\begin{ex} %[4H6V30]
Qua một điểm nằm ngoài một đường thẳng, vẽ được bao nhiêu đường thẳng song song với đường thẳng đã cho?
\shortans{1}
\loigiai{
\textbf{Đáp số:} $1$\\
1 đường thẳng duy nhất.
}
\end{ex}

\begin{ex} %[4H6V30]
Hình vuông có bao nhiêu cặp cạnh đối diện song song?
\shortans{2}
\loigiai{
\textbf{Đáp số:} $2$\\
2 cặp.
}
\end{ex}

\begin{ex} %[4H6V30]
Một khung cửa sổ hình chữ nhật có các chấn song sắt đặt cách đều nhau. Các chấn song này có quan hệ gì với nhau (song song hay vuông góc)? Ghi "song song".
\shortans{song song}
\loigiai{
\textbf{Đáp số:} song song\\
Các chấn song sắt song song với nhau.
}
\end{ex}

\begin{ex} %[4H6V30]
Nếu $a // b$ và $b // c$ thì $a$ và $c$ có quan hệ gì với nhau? Ghi "song song".
\shortans{song song}
\loigiai{
\textbf{Đáp số:} song song\\
$a // c$.
}
\end{ex}

\begin{ex} %[4H6V30]
Một hình thang cân có mấy cạnh đáy song song với nhau?
\shortans{2}
\loigiai{
\textbf{Đáp số:} $2$\\
2 cạnh đáy song song.
}
\end{ex}

\begin{ex} %[4H6V30]
Cho 4 đường thẳng song song phân biệt. Chúng cắt một đường thẳng thứ năm tại 4 điểm. Có bao nhiêu đoạn thẳng được tạo thành trên đường thẳng thứ năm bởi 4 điểm đó?
\shortans{6}
\loigiai{
\textbf{Đáp số:} $6$\\
$4 \times 3 : 2 = 6$ đoạn thẳng.
}
\end{ex}
